% !TeX root = ../../Book3.tex
% 纲要标题：### 18.3 局部化与完备化
% 纲要位置：计划纲要.md，第 2064 行。

\section{局部化与完备化}
\label{b3:sec:1803}

% 纲要内容：
%
% * 正则性在局部化下的行为
% * 正则性与完备化
% * 正则局部环是整环的证明方法
%

\subsection{正则性的局部化}
\label{b3:subsec:1803:01}

\begin{theorem}\label{b3:thm:regular-local-localization}
设 \(R\) 为交换 Noether 正则局部环，\(\mathfrak p\in\operatorname{Spec}R\)。则 \(R_{\mathfrak p}\) 仍是正则局部环。
\end{theorem}
\begin{proof}
有限 \(R\) 模 \(R/\mathfrak p\) 由\cref{b3:thm:regular-local-global-dimension}具有有限自由消解。局部化正合，故该消解在 \(\mathfrak p\) 处给出
\((R/\mathfrak p)_{\mathfrak p}=\kappa(\mathfrak p)\) 的有限自由 \(R_{\mathfrak p}\) 消解。后者是 Noether 局部环的剩余域，再用同一定理的反向便得正则性。
\end{proof}

\begin{definition}\label{b3:def:regular-ring}
设 \(R\) 为交换 Noether 环。若每个素理想 \(\mathfrak p\) 的局部环 \(R_{\mathfrak p}\) 都正则，则称 \(R\) 为\term[zhengzehuan]{正则环}[regular ring]。
\end{definition}

只检验极大理想已足够：每个素理想包含在某个极大理想中，再使用前一定理。正则环未必局部，也未必整环，例如两个域的直积。关于正则局部环的断言若去掉“局部”，需要重新检查结论。

\subsection{正则性与完备化}
\label{b3:subsec:1803:02}

\begin{theorem}\label{b3:thm:regular-local-completion}
设 \((R,\mathfrak m,k)\) 为交换 Noether 局部环，\(\widehat R\) 为其 \(\mathfrak m\) 进完备化。则 \(R\) 正则当且仅当 \(\widehat R\) 正则。
\end{theorem}
\begin{proof}
由\cref{b3:thm:completion-noetherian}，\(\widehat R\) 是 Noether 局部环，极大理想为 \(\mathfrak m\widehat R\)，剩余域仍为 \(k\)。\cref{b3:prop:completion-quotients}在二次幂处给出
\[
\mathfrak m/\mathfrak m^2
\cong\mathfrak m\widehat R/(\mathfrak m\widehat R)^2,
\]
所以两环嵌入维数相同。\cref{b3:thm:completion-dimension}使它们 Krull 维数相同。定义中的两个数分别不变，故正则性等价。
\end{proof}

例如 \(k[x_1,\ldots,x_d]_{(x_1,\ldots,x_d)}\) 的完备化为 \(k[[x_1,\ldots,x_d]]\)，所以后者正则。这里保留全部形式幂级数，不能把它与只允许有限项的多项式局部环当成同一个环。

\subsection{正则局部环的整性}
\label{b3:subsec:1803:03}

\begin{proposition}\label{b3:prop:regular-local-domain}
交换 Noether 正则局部环是整环。
\end{proposition}
\begin{proof}
设极大理想为 \(\mathfrak m\)。若 \(a,b\ne0\)，Krull 交定理使它们具有有限阶数 \(r,s\)，故在伴随分次环中给出非零初始式。\cref{b3:thm:regular-local-characterizations}说明该分次环是域上的多项式环，因而两初始式的乘积非零。这一乘积正是 \(ab\) 在 \(\mathfrak m^{r+s}/\mathfrak m^{r+s+1}\) 中的像，所以 \(ab\ne0\)。
\end{proof}

这个证明也说明分离性和分次结构各自的作用：前者保证非零元素有初始式，后者保证相乘时最低项不会消失。一般环的伴随分次环不一定是整环，即使原环是整环；尖点的切锥便给出反例。

\ProseParagraphBreak
一维正则局部环的极大理想由一个元素生成，再结合刚得到的整性，它由\cref{b3:thm:dvr-characterizations}为 DVR。下一节将把唯一分解推广到任意有限维数；这需要比“一维局部化是 DVR”更多的信息。
